\chapter{Typed appendix and near-$126$ computation data}
\label{chap:appendix-data}

The paper's appendix contains 401 nonempty class rows and nine explicit empty
filtration bands.  All 410 semantic records are in scope.  The intended Lean
encoding is finite, typed data carried by external-evidence values; it is not a
collection of project axioms.  Every duplicated display of a differential
shares one relation identifier.

\section{Data schema and provenance}

% SOURCE: aimpaper/main.tex:2738-3290, all twelve tables.
\begin{definition}[Appendix table identifier]
  \label{def:appendix-table-id}
  \notready
  The table identifier is a twelve-element finite type containing
  $C\nu126$, $S122$, $S123$, the two $S124$ bands, the two $S125$ bands, the
  two $S126$ bands, and the three $S127$ bands.  It records the paper table
  number, spectrum, stem, filtration interval, source line range, and PDF page.
\end{definition}

% SOURCE: aimpaper/main.tex:2738-3290, class-expression column; aimpaper/112.tex:8-1026.
\begin{definition}[Typed Ext class expression]
  \label{def:typed-class-expression}
  \uses{def:steenrod-ext,def:hopf-cofiber-cell-notation}
  \notready
  A typed class expression is a sparse $\mathbb F_2$ expression built from a
  dataset-versioned generator $x_{n,s,k}$, a named class, product, and finite
  sum, together with spectrum, stem, filtration, internal degree, and optional
  bottom- or top-cell tag.  Constructors prove degree compatibility; a finite
  candidate set represents ambiguity without converting it to free-form text.
\end{definition}

% SOURCE: aimpaper/main.tex:2738-3290, differential/status columns.
\begin{definition}[Appendix row status]
  \label{def:appendix-row-status}
  \uses{def:typed-class-expression}
  \notready
  Row status is one of:
  \begin{enumerate}
    \item an explicit zero filtration band;
    \item a permanent class;
    \item a known outgoing $d_r$ with a nonempty finite target set;
    \item a known incoming $d_r$ with a nonempty finite source set;
    \item survival through a stated page with the next outgoing differential
      unresolved.
  \end{enumerate}
  An incoming marker $d_r^{-1}$ is not an inverse map.  A printed $d_r?$ does
  not assert existence of a nonzero differential.  A ``possibly'' term is a
  finite target ambiguity attached to a known relation.
\end{definition}

% SOURCE: aimpaper/main.tex:2738-3290; duplicated source/target presentations across adjacent stems.
\begin{definition}[Differential relation identifier]
  \label{def:differential-relation-id}
  \uses{def:appendix-row-status}
  \notready
  Each known differential has a stable relation identifier, length, typed
  source expression, and typed target candidate set.  The outgoing row in one
  stem and incoming row in the next refer to the same identifier.  A validation
  rule rejects two records with incompatible length, source, or target data.
\end{definition}

% SOURCE: PROJECT_BOUNDARY.md, Appendix data; aimpaper/main.tex:2785-2800.
\begin{definition}[Appendix evidence record]
  \label{def:appendix-evidence-record}
  \uses{def:appendix-table-id,def:appendix-row-status,def:differential-relation-id,def:external-evidence}
  \notready
  An appendix evidence record contains a table identifier, a unique row key,
  spectrum and bidegree, typed class expression or zero band, row status,
  relation identifier when applicable, paper locator, computation-artifact
  locator, program version, input hash, method, and an explicit evidence proof
  of the proposition represented by that status.
\end{definition}

% SOURCE: aimpaper/main.tex:213-240 and 2785-2800; LWXMachine and LWXZenodo bibliography entries.
\begin{proposition}[Lin computation provenance]
  \label{prop:lin-computation-provenance}
  \uses{def:appendix-evidence-record,def:source-ref}
  \notready
  The near-$126$ dataset is tied to the Lin program, Zenodo record, input
  spectra and maps, program commit or archive hash, output artifact, and proof
  identifiers for propagated or disproved differentials.  Until every field
  is populated, the dataset remains not ready; the interactive plot alone is
  not sufficient provenance.
\end{proposition}

\begin{proof}
  Cross-check the paper bibliography against the local source inventory, open
  the archived computation files, and record stable paths and hashes.  Verify
  every row key against the artifact that proves or bounds its status.  Reject
  records justified only by visual coordinates or an unversioned generator
  name.
\end{proof}

% SOURCE: aimpaper/main.tex:2785-2792, three manually supplied differentials.
\begin{proposition}[Manual differential inputs]
  \label{prop:appendix-manual-inputs}
  \uses{prop:lin-computation-provenance,def:external-result,def:external-evidence}
  \notready
  The computation input ledger separately records
  \[
  \begin{aligned}
    d_5(h_0^{24}h_6)&=h_0^2P^6d_0,\\
    d_6(h_0^{55}h_7)&=h_0^2x_{126,60},\\
    d_3(v_2^{16})&=\beta^5g\quad\text{in }tmf.
  \end{aligned}
  \]
  Each is an explicit external result or evidence input with a theorem, page,
  or archived computation locator.  The two image-of-$J$ locators and the
  precise Bruner--Rognes locator are currently provenance obligations.
\end{proposition}

\begin{proof}
  Type-check the three source and target degrees, attach their independent
  source locators, and verify that the archived program declares exactly these
  inputs.  Do not allow a table output derived from an input to serve as the
  source of that same input.
\end{proof}

% SOURCE: aimpaper/main.tex:2796-2800 and repeated rows at lines 2999, 3096, 3161.
\begin{proposition}[$x_{126,21}$ finite target ambiguity]
  \label{prop:x12621-ambiguity}
  \uses{def:appendix-evidence-record,prop:lin-computation-provenance}
  \notready
  The known $d_4$ from $x_{126,21}$ has target candidate
  \[
    x_{125,25,2}+x_{125,25}+g^4\Delta h_1g
      +\epsilon d_0^2e_0gB_4,
    \qquad \epsilon\in\mathbb F_2,
  \]
  rather than a uniquely determined target.  All repeated table displays
  refer to one relation identifier with this two-element candidate set.
\end{proposition}

\begin{proof}
  Transcribe the prose and the three relevant table rows, normalize their class
  expressions, and verify that their fixed part agrees.  Encode the optional
  summand as the finite coefficient $\epsilon$, then link every display to the
  shared relation identifier.
\end{proof}

\section{The twelve table datasets}

% SOURCE: aimpaper/main.tex:2738-2773, Table Table:Cnu126, PDF p.54.
\begin{proposition}[$S^0/\nu$ stem-$126$ table]
  \label{data:table-cnu126}
  \uses{def:appendix-evidence-record,prop:lin-computation-provenance}
  \notready
  Table 1 contains exactly 26 nonempty rows in stem $126$ and filtrations
  $9\le s\le14$: 12 known incoming, 10 known outgoing, one permanent, and
  three unresolved rows, with no zero band.
\end{proposition}

\begin{proof}
  Parse lines 2738--2773, normalize the cell tags and expressions, and prove
  by finite enumeration that the four status counts sum to 26 and every
  printed class row occurs exactly once.
\end{proof}

% SOURCE: aimpaper/main.tex:2804-2848, Table Table:S122, PDF p.56.
\begin{proposition}[$S^0$ stem-$122$ table]
  \label{data:table-s122}
  \uses{def:appendix-evidence-record,prop:lin-computation-provenance}
  \notready
  Table 2 contains 35 nonempty rows: 16 incoming, 13 outgoing, five permanent,
  and one unresolved.  It also contains the zero bands $s=9$--$10$ and
  $s=0$--$7$.
\end{proposition}

\begin{proof}
  Transcribe lines 2804--2848 and enumerate the status constructors and two
  zero-band records.  Check the table's filtration coverage through $s=25$.
\end{proof}

% SOURCE: aimpaper/main.tex:2851-2903, Table Table:S123, PDF p.57.
\begin{proposition}[$S^0$ stem-$123$ table]
  \label{data:table-s123}
  \uses{def:appendix-evidence-record,prop:lin-computation-provenance}
  \notready
  Table 3 contains 42 nonempty rows: 22 incoming, 17 outgoing, two permanent,
  and one unresolved.  Its zero bands are $s=25$, $s=21$--$22$, and
  $s=0$--$7$.
\end{proposition}

\begin{proof}
  Transcribe lines 2851--2903, verify all row bidegrees and the three zero
  bands, and prove the status counts by finite enumeration.
\end{proof}

% SOURCE: aimpaper/main.tex:2908-2958, Table Table:S124.13, PDF p.58.
\begin{proposition}[$S^0$ stem-$124$ high-filtration table]
  \label{data:table-s124-high}
  \uses{def:appendix-evidence-record,prop:lin-computation-provenance}
  \notready
  Table 4 contains 43 nonempty rows for $13\le s\le25$: 24 incoming, 13
  outgoing, and six permanent, with no unresolved row or zero band.
\end{proposition}

\begin{proof}
  Transcribe lines 2908--2958 and verify by enumeration that the three status
  counts sum to 43 and cover every printed filtration in the stated band.
\end{proof}

% SOURCE: aimpaper/main.tex:2962-2992, Table Table:S124.12, PDF p.59.
\begin{proposition}[$S^0$ stem-$124$ low-filtration table]
  \label{data:table-s124-low}
  \uses{def:appendix-evidence-record,prop:lin-computation-provenance}
  \notready
  Table 5 contains 23 nonempty rows for $s\le12$: 11 incoming, 11 outgoing,
  and one permanent.  It contains the zero band $s=0$--$5$.
\end{proposition}

\begin{proof}
  Transcribe lines 2962--2992, check the zero band and all class degrees, and
  prove the three status counts by finite enumeration.
\end{proof}

% SOURCE: aimpaper/main.tex:2994-3018, Table Table:S125.20, PDF p.59.
\begin{proposition}[$S^0$ stem-$125$ high-filtration table]
  \label{data:table-s125-high}
  \uses{def:appendix-evidence-record,prop:lin-computation-provenance,prop:x12621-ambiguity}
  \notready
  Table 6 contains 17 nonempty rows for $20\le s\le25$: six incoming, ten
  outgoing, and one permanent.  One outgoing row carries the finite
  ``possibly'' ambiguity and there is no zero band.
\end{proposition}

\begin{proof}
  Transcribe lines 2994--3018, connect the ambiguous row to the shared relation
  identifier, and enumerate the remaining statuses and bidegrees.
\end{proof}

% SOURCE: aimpaper/main.tex:3020-3078, Table Table:S125.19, PDF p.60.
\begin{proposition}[$S^0$ stem-$125$ low-filtration table]
  \label{data:table-s125-low}
  \uses{def:appendix-evidence-record,prop:lin-computation-provenance}
  \notready
  Table 7 contains 50 nonempty rows for $s\le19$: 18 incoming, 28 outgoing,
  two permanent, and two unresolved.  It contains the zero band
  $s=0$--$4$.
\end{proposition}

\begin{proof}
  Transcribe lines 3020--3078 and prove the counts and filtration coverage.
  Preserve the zero band as a proposition-valued record because it is used in
  the main reduction to exclude low-weight $\lambda$-torsion.
\end{proof}

% SOURCE: aimpaper/main.tex:3081-3135, Table Table:S126.11, PDF p.61.
\begin{proposition}[$S^0$ stem-$126$ high-filtration table]
  \label{data:table-s126-high}
  \uses{def:appendix-evidence-record,prop:lin-computation-provenance,prop:x12621-ambiguity}
  \notready
  Table 8 contains 47 nonempty rows for $11\le s\le25$: 29 incoming, 15
  outgoing, one permanent, and two unresolved.  One row has finite target
  ambiguity and there is no zero band.
\end{proposition}

\begin{proof}
  Transcribe lines 3081--3135, link the ambiguous display, and verify all four
  counts and the complete filtration interval by finite enumeration.
\end{proof}

% SOURCE: aimpaper/main.tex:3138-3171, Table Table:S126.10, PDF p.62.
\begin{proposition}[$S^0$ stem-$126$ low-filtration table]
  \label{data:table-s126-low}
  \uses{def:appendix-evidence-record,prop:lin-computation-provenance,prop:x12621-ambiguity}
  \notready
  Table 9 contains 26 nonempty rows for $s\le10$: 11 incoming, ten outgoing,
  one permanent, and four unresolved.  One row has finite ambiguity; the zero
  band is $s=0$--$1$.  In particular the $h_6^2$ row is unresolved, not a
  known nonzero $d_7$.
\end{proposition}

\begin{proof}
  Transcribe lines 3138--3171, preserve each question-mark status, link the
  ambiguity relation, and verify counts and the single zero band.
\end{proof}

% SOURCE: aimpaper/main.tex:3173-3197, Table Table:S127.21, PDF p.62.
\begin{proposition}[$S^0$ stem-$127$ high-filtration table]
  \label{data:table-s127-high}
  \uses{def:appendix-evidence-record,prop:lin-computation-provenance}
  \notready
  Table 10 contains 17 nonempty rows for $21\le s\le25$: five incoming, 11
  outgoing, and one permanent, with no unresolved row or zero band.
\end{proposition}

\begin{proof}
  Transcribe lines 3173--3197 and verify all class degrees, relation links, and
  status counts by finite enumeration.
\end{proof}

% SOURCE: aimpaper/main.tex:3200-3256, Table Table:S127.20, PDF p.63.
\begin{proposition}[$S^0$ stem-$127$ middle-filtration table]
  \label{data:table-s127-middle}
  \uses{def:appendix-evidence-record,prop:lin-computation-provenance}
  \notready
  Table 11 contains 50 nonempty rows for $10\le s\le20$: 20 incoming, 22
  outgoing, and eight permanent, with no unresolved row or zero band.
\end{proposition}

\begin{proof}
  Transcribe lines 3200--3256 and prove by finite enumeration that all 50 rows
  have one of the three stated statuses and consistent relation identifiers.
\end{proof}

% SOURCE: aimpaper/main.tex:3258-3290, Table Table:S127.9, PDF p.64.
\begin{proposition}[$S^0$ stem-$127$ low-filtration table]
  \label{data:table-s127-low}
  \uses{def:appendix-evidence-record,prop:lin-computation-provenance}
  \notready
  Table 12 contains 25 nonempty rows for $s\le9$: four incoming, 14 outgoing,
  two permanent, and five unresolved.  Its zero band is $s=0$.
\end{proposition}

\begin{proof}
  Transcribe lines 3258--3290, preserve all five unresolved statuses, and
  verify the counts, bidegrees, and zero band by finite enumeration.
\end{proof}

\section{Completeness and the near-$126$ visualization}

% SOURCE: aimpaper/main.tex:2738-3290, all tables; local PDF pp.54,56-64.
\begin{theorem}[All appendix rows are encoded]
  \label{thm:appendix-all-rows-encoded}
  \uses{data:table-cnu126,data:table-s122,data:table-s123,data:table-s124-high,data:table-s124-low,data:table-s125-high,data:table-s125-low,data:table-s126-high,data:table-s126-low,data:table-s127-high,data:table-s127-middle,data:table-s127-low,prop:appendix-manual-inputs}
  \notready
  The twelve datasets contain exactly 401 nonempty class records and nine zero
  bands.  Among the nonempty records there are 178 incoming, 174 known
  outgoing, 31 permanent, and 18 unresolved statuses.  Every printed row and
  zero band appears exactly once, every duplicated differential is linked by
  one relation identifier, and no unprinted row is introduced.
\end{theorem}

\begin{proof}
  Concatenate the twelve finite tables and evaluate decidable checks for key
  uniqueness, source-line coverage, spectrum and bidegree consistency, status
  totals, and paired relation identifiers.  Compare the normalized
  serialization with the fixed source hashes and emit a coverage certificate
  listing every source row and zero band.
\end{proof}

% SOURCE: aimpaper/112.tex:5-1029, Figure Fig:near126; included at aimpaper/main.tex:2152; local PDF p.43.
\begin{definition}[Near-$126$ figure contract]
  \label{def:near126-figure-contract}
  \uses{thm:appendix-all-rows-encoded}
  \notready
  Figure 9 is a derived visualization over stems $122$--$127$ and filtrations
  $0$--$19$.  Its left panel has 289 dots, 384 black structural edges, and 157
  colored known differentials (87 $d_2$, 44 $d_3$, 17 $d_4$, seven $d_5$,
  and two $d_7$).  Its right panel has 42 dots, 21 black edges, and 12 dashed
  shortest-potential differentials.  Dashed arrows are unresolved candidates,
  and the figure is not accepted as independent evidence because its source
  has no complete legend and uses layout offsets for multiple classes.  Lines
  1031--1919 of \texttt{112.tex} are commented duplicates and are ignored.
\end{definition}

% VERIFIED-IN: local source hashes recorded during audit.
\begin{proposition}[Primary-source snapshot]
  \label{prop:paper-source-snapshot}
  \uses{def:source-ref}
  \notready
  The audited snapshot records 66 PDF pages and SHA256 values
  \texttt{7cae269851a88d10dd194651dbf7497b75ef8b8b914bc1901dfa3734cd8096b4}
  for the PDF,
  \texttt{1125462bcae4a4ec56e3bfcaad15df4febf98757dfb83462b155af162c99c9e0}
  for \texttt{main.tex}, and
  \texttt{5eb83ef105ce9dbbda87a9f57726bb5c078308fd3718840ea29ee8100cbefabb}
  for \texttt{112.tex}.
\end{proposition}

\begin{proof}
  Recompute the three hashes, verify the PDF page count, and compare them with
  the stored source-reference fields before accepting any row-completeness
  certificate.
\end{proof}
